%%%%%%%%%%  scan idx 62  |  folio 49  %%%%%%%%%%
du début du \S 2, on
définit une multiplication dans $\underline{K}(X)$ en posant
\begin{equation}\tag{2.15}
 \gamma(F)\gamma(G)=\sum_{i\geq0}(-1)^i\gamma(\underline{\Tor}_i(F,G))
\end{equation}
(les $\underline{\Tor}_i$ étant pris, bien entendu, par rapport aux anneaux locaux de $X$). La suite exacte des $\underline{\Tor}_i$ montre que la formule (2.15) définit bien une application biadditive de $\underline{K}(X)\times\underline{K}(X)$ dans $\underline{K}(X)$~; l'associativité résulte de la suite spectrale des $\Tor$ multiples. On vérifie aussi de cette façon que l'on peut calculer le produit de plusieurs facteurs par la formule~:
\begin{equation}\tag{2.16}
 \gamma(F_1)....\gamma(F_p)=\sum_{i\geq0}(-1)^i\gamma(\underline{\Tor}_i(F_1,\ldots,F_p))
\end{equation}
où, dans le deuxième membre, les $\underline{\Tor}_i$ sont des «~$\underline{\Tor}$ simultanés~», calculés par résolution projective simultanée des arguments $F_i$.

Si l'on suppose dans (2.15) que $G$ est \emph{localement libre}, on trouve~:
\begin{equation}\tag{2.15 bis}
 \gamma(F)\gamma(G)=\gamma(F\otimes G),
\end{equation}
ce qui montre en particulier que $\gamma(\cO_X)$ est élément unité de $\underline{K}(X)$, et que l'homomorphisme canonique de $\underline{K}(\xi(X))$ dans $\underline{K}(X)$ est un homomorphisme d'anneaux.

La formule (2.15) est évidemment suggérée par la «~formule des $\Tor$~» de Serre en théorie des intersections. Cette formule prouve la deuxième partie de la proposition suivante, la première résultant immédiatement de la formule (2.15) et d'un calcul de $\underline{\Tor}_i$ locaux~:

\textbf{Proposition 2.7.} \emph{Soient $X$ une variété non singulière, $Y$ et $Z$ deux cycles dans $X$. Si $Y$ et $Z$ se coupent partout transversalement, on a~:}
\begin{equation}\tag{2.17}
 \gamma(Y)\gamma(Z)=\gamma(Y.Z).
\end{equation}
\emph{Si $Y$ et $Z$ sont homogènes et de dimension respective $n-i$ et $n-j$, et s'ils se coupent sans composante excédentaire, on a alors~:}
